(1) 腔内光单次往返光程是 \(2nL\)，故频率 \(\nu\) 对应两次相邻透射光的传播相位差 \(\delta=4\pi nL/\lambda_0=4\pi nL\nu/c\)。在两界面相位可看作固定的对称无损镜近似下，相邻透射复振幅比为 \(Re^{i\delta}\)。
(2) 将各次透射振幅作几何级数相加，得到 \(E_t/E_i=(1-R)e^{i\phi}/(1-Re^{i\delta})\)，其中公共相位 \(\phi\) 不影响强度。取模平方并使用 \(|1-Re^{i\delta}|^2=(1-R)^2+4R\sin^2(\delta/2)\)，得 \(\mathcal T=1/[1+F\sin^2(\delta/2)]\)，\(F=4R/(1-R)^2\)。
(3) 谐振峰满足 \(\delta=2\pi m\)，故相邻峰频率间隔 \(\mathrm{FSR}=c/(2nL)\)。在峰附近 \(\sin(\delta/2)\simeq\Delta\delta/2\)，半最大值要求 \(F(\Delta\delta/2)^2=1\)，全宽为 \(\delta_{\rm FWHM}=4/\sqrt F\simeq2(1-R)/\sqrt R\)。精细度 \(\mathcal F=2\pi/\delta_{\rm FWHM}\simeq\pi\sqrt R/(1-R)\)。
(4) \(R\to0\) 时多次往返消失，透射接近一且谐振峰不明显，高反射率线宽近似失效；\(R\to1^-\) 时峰值仍为一而线宽趋窄。公式中透射镜损耗为零，若存在吸收与散射损失，谐振峰最大透射一般低于一，不能使用相同的归一公式。
